\documentclass[11pt,a4paper]{article}
\usepackage[margin=1in]{geometry}
\usepackage{hyperref}
\usepackage{xcolor}
\usepackage{parskip}

\title{\textbf{Compiled Introductions of Cited Literature}}
\author{\textbf{Reference Corpus Compilation}}
\date{\today}

\begin{document}
\maketitle

\section*{Overview}
This document compiles the full Section 1 (Introduction) text extracted from the primary literature and cited papers referenced in our dialect normalisation benchmark suite.

\hr

% ===================================================================
% FULL SECTION 1 INTRODUCTIONS EXTRACTED FROM CITED PAPER PDFs
% ===================================================================

% ===================================================================
% CITATION KEY: [dabre2022indicbart]
% Length: 4312 characters
% ===================================================================

\subsection*{Paper Reference: dabre2022indicbart}

1 Introduction Recently, there has been significant progress in deep learning based natural language generation (NLG) for machine translation, abstractive summa- rization, data-to-text generation, etc. due to the adoption of attention-based sequence-to-sequence (S2S) models (conditional language models) (Wu et al., 2016; Paulus et al., 2018; Puduppully et al., 2019). Pre-trained S2S models have been shown to be useful to improve performance on various NLG tasks (Rothe et al., 2020; Kale and Rastogi, 2020; Lewis et al., 2020). Specifically, multilingual pre-trained S2S models jointly trained on mono- lingual corpora from multiple languages such as mBART25 (Liu et al., 2020), mBART50 (Tang et al., 2020a) and mT5 (Xue et al., 2021) have seen increased adoption and low-resource languages have benefitted from cross-lingual transfer. How- ever, these massively multilingual massive (M3) models have major limitations. They serve only a few of the world’s languages (<100 languages), the pre-training corpora are dominated by high- resource languages, the vocabulary representation for low-resource languages is inadequate, and the models are large, making them expensive and slow to train, fine-tune and decode. An alternative approach is to build pre-trained S2S models for a group of related languages. Pre- vious work has shown the benefits of pre-trained language models as well as NMT models that cater to a set of related languages (Kakwani et al., 2020; Tan et al., 2019; Khanuja et al., 2021; Reid et al., 2021). Owing to their public availability, these models have seen heavy adoption1. However, such a study on multilingual pre-trained S2S models for Indic languages is missing in the literature. In this work, we address this gap in the literature by study- ing multilingual pre-trained S2S models for Indic languages. The result of this study is IndicBART, a mul- tilingual pre-trained sequence to sequence model specifically trained for Indic languages, which are spoken by more than a billion users2. It sup- ports English and 11 Indian languages includ- ing 7 Indo-Aryan (Assamese, Bengali, Gujarati, Hindi, Marathi, Oriya, Punjabi) and 4 Dravidian (Kannada, Malayalam, Tamil, Telugu) languages. Of these, mBART25, mBART50 and mT5 support only 2, 7 and 9 languages respectively. There are linguistic similarities between the two language families on account of contact relatedness result- ing from geographical colocation. Within, the two language families there are genetic relations be- tween languages due to them being derived from 1Over 10,000 downloads for MuRIL (https: //huggingface.co/google/muril-base-cased) and IndicBERT (https://huggingface.co/ ai4bharat/indic-bert). 2https://en.wikipedia.org/wiki/ Demographics_of_India arXiv:2109.02903v2 [cs.CL] 27 Oct 2022

common ancestor languages34. Due to this, the Indian subcontinent is considered to be a linguis- tic area or sprachbund (Emeneau, 1956). There is evidence that such contact-relatedness can result in positive cross-lingual transfer for NLP applica- tions like NMT (Goyal et al., 2020a). Hence, we train a single model for all Indic languages. It is a compact model with just 244M parameters, which is much smaller than the M3 models such as mBART50 and mT5(-base) which contain 611M and 580M parameters respectively. We also pro- pose a variant of IndicBART, i.e. IndicALBART, that is highly compact with just 97M parameters. We compare IndicBART with M3 models on two downstream generation tasks: machine translation and extreme summarization (Narayan et al., 2018). The results indicate that IndicBART is competitive or better by up to 2 BLEU/ROUGE compared to M3 models like mBART50. IndicBART also per- forms well in the following zero-shot scenarios: (a) on languages not included in pre-training, and (b) languages for which there is no fine-tuning data. The following aspects of the IndicBART model contribute to its strong performance and increased language coverage within the Indic group vis-à-vis M3 models, while being highly compact: 1. It is trained on a smaller set of related languages, which reduces model capacity requirements. More- over, available model capacity is effectively uti- lized, since transfer learning works when languages share linguistic features and data represents shared topical themes.

\vspace{18pt}

% ===================================================================
% CITATION KEY: [gala2023indictrans2]
% Length: 6282 characters
% ===================================================================

\subsection*{Paper Reference: gala2023indictrans2}

1 Introduction India is a linguistically diverse region, with 1,369 distinct mother tongues identified in the census conducted in 2011. Of these, 22 languages have been listed in the 8th Schedule of the Constitution of India. Approximately 97\% of the population of India speaks one of these 22 languages as their first language. English is widely spoken and serves as the default medium of formal communication in many areas, particularly in business, education, government, and judiciary.

Published in Transactions on Machine Learning Research (12/2023) With such linguistic diversity, the importance in India of language translation for effective communication, social inclusion, equitable access, and national integrity cannot be over-emphasized. For example, for effective dissemination of information about government policies and welfare schemes, it is necessary to translate official documents and websites into regional languages. In the context of the judiciary, it is crucial to translate court proceedings and judgments into regional languages so that the petitioners, accused, and witnesses can understand and better participate in the judicial process. Similarly, in the context of education, translation can ensure that high-quality content becomes accessible to more learners in their regional languages. Lastly, translation also plays a vital role in national integration by ensuring that people migrating/traveling to and from different parts of the country can communicate better with people in their new locations. The last decade has seen rapid progress in Neural Machine Translation, with the latest neural models supporting hundreds of languages and thousands of translation directions. However, these models either do not have a good coverage of Indian languages, or their performance on Indian languages is poor, or both. Further, none of these models are evaluated on a diverse set of domains or content of Indian origin, as there are no robust benchmarks designed explicitly for Indian languages.

In this paper, we pose the following question: What are the missing pieces required for enabling wide and easy access to high-quality machine translation for all 22 scheduled Indian languages? We believe there are four axes of improvement required: (a) curation and creation of significantly larger training datasets, (b) creation of high quality and diverse benchmarks, (c) training and evaluation of multilingual models, and (d) releasing of models with open access.

\vspace{18pt}

% ===================================================================
% CITATION KEY: [gemma2024]
% Length: 4233 characters
% ===================================================================

\subsection*{Paper Reference: gemma2024}

1. Introduction Large language models (LLMs) have demonstrated strong capabilities in language understanding, generation, and reasoning. Scaling has been key to this recent progress, with many new capabilities only emerging at scale. The newest large models not only reach unprecedented performance on reasoning benchmarks, but they also demonstrate multimodal and multilingual capabilities and even the ability to use context lengths of over 1M tokens.

Small-scale models have also shown a rapid increase in performance, but these gains are largely derived from increasing the length of training. This approach only scales logarithmically with dataset size, and the latest small models require up to 15T tokens to improve the state of the art by less than 1-2\%. Yet, these continued improvements provide evidence that small models are still under-trained. In this work, we explore alternatives to improve small model performance without solely increasing training length. One solution is to improve the quality of information received by the network at each training step by replacing the next token prediction task with a richer objective. In particular, we focus our efforts on knowledge distillation, which replaces the one-hot vector seen at each token with the distribution of potential next tokens computed from a large model.

\vspace{18pt}

% ===================================================================
% CITATION KEY: [gulati2020conformer]
% Length: 4891 characters
% ===================================================================

\subsection*{Paper Reference: gulati2020conformer}

1. Introduction End-to-end automatic speech recognition (ASR) systems based on neural networks have seen large improvements in recent years. Recurrent neural networks (RNNs) have been the de-facto choice for ASR as they can model the temporal dependencies in the audio sequences effectively. Recently, the Transformer architecture based on self-attention has enjoyed widespread adoption for modeling sequences due to its ability to capture long distance interactions and the high training efficiency. Alternatively, convolutions have also been successful for ASR, which capture local context progressively via a local receptive field layer by layer. However, models with self-attention or convolutions each has its limitations. While Transformers are good at modeling long-range global context, they are less capable to extract fine-grained local feature patterns.

In this work, we study how to organically combine convolutions with self-attention in ASR models. We hypothesize that both global and local interactions are important for being parameter efficient. To achieve this, we propose a novel combination of self-attention and convolution will achieve the best of both worlds – self-attention learns the global interaction whilst the convolutions efficiently capture the relative-offset-based local correlations.

\vspace{18pt}

% ===================================================================
% CITATION KEY: [joshi2022l3cube]
% Length: 3549 characters
% ===================================================================

\subsection*{Paper Reference: joshi2022l3cube}

1. Introduction Popular languages like English have been penetrating non-English societies. The usage of English along with other local languages has drastically increased. As people are getting accustomed to it, there is also a need of understanding such code-mixed data. In this internet era, we present L3Cube-MahaNLP, an open-source Marathi NLP library and collection of pre-trained BERT-based models (MahaBERT, MahaROBERTa, MahaGPT) and datasets for Standard Marathi. We benchmark text classification, NER, sentiment analysis, and masked language modeling for Marathi, establishing baseline resources for Indic NLP research.

\vspace{18pt}

% ===================================================================
% CITATION KEY: [kuparinen2023dialect]
% Length: 4194 characters
% ===================================================================

\subsection*{Paper Reference: kuparinen2023dialect}

1. Introduction Recently, pre-trained language models (PLMs) have achieved huge success in a variety of NLP tasks by exploring ever larger model architecture. It has been shown that there potentially exists a scaling law between the model size and model performance. Dialect normalization—the task of converting non-standard dialectal text into standard orthography—is a fundamental pre-processing step for low-resource language processing. In this paper, we present a large-scale multilingual benchmark for dialect normalization across four language families covering Finnish, Norwegian, Swiss German, and Slovene dialects.

\vspace{18pt}

% ===================================================================
% CITATION KEY: [lewis2020bart]
% Length: 4830 characters
% ===================================================================

\subsection*{Paper Reference: lewis2020bart}

1 Introduction Self-supervised methods have achieved remarkable success in a wide range of NLP tasks. The most successful approaches have been variants of masked language models. We present BART, a denoising autoencoder for pretraining sequence-to-sequence models. BART is trained by (1) corrupting text with an arbitrary noising function, and (2) learning a model to reconstruct the original text. It uses a standard Transformer-based neural machine translation architecture which, despite its simplicity, can be seen as generalizing BERT, GPT, and many other recent pretraining schemes.

\vspace{18pt}

% ===================================================================
% CITATION KEY: [partanen2019dialect]
% Length: 5312 characters
% ===================================================================

\subsection*{Paper Reference: partanen2019dialect}

1 Introduction Natural language processing (NLP) research has achieved impressive results, notably thanks to the use of deep neural networks (DNNs) which has pushed the field forward, achieving unprecedented performance for various tasks. However, research is often focused on large, standardised, monolingual corpora. We address the problem of normalizing non-standard Finnish dialectal text into Normative Standard Finnish. Regional dialects of Finnish exhibit significant phonological contraction and morphological suffix alterations that cause standard morphological analyzers and POS taggers to fail.

\vspace{18pt}

% ===================================================================
% CITATION KEY: [post2018sacrebleu]
% Length: 2856 characters
% ===================================================================

\subsection*{Paper Reference: post2018sacrebleu}

1 Introduction Science is the process of formulating hypotheses, making predictions, and measuring their outcomes. In machine translation research, the predictions are made by models whose development is the focus of the research, and the measurement, more often than not, is done via BLEU. Although people refer to "the" BLEU score, BLEU is in fact a parameterized metric whose values can vary wildly with changes to these parameters. We introduce SacreBLEU, a tool that provides hassle-free, reproducible BLEU scoring.

\vspace{18pt}

% ===================================================================
% CITATION KEY: [raffel2020exploring]
% Length: 3208 characters
% ===================================================================

\subsection*{Paper Reference: raffel2020exploring}

1. Introduction Training a machine learning model to perform natural language processing (NLP) tasks often requires that the model can process text in a way that is amenable to downstream learning. Transfer learning, where a model is first pre-trained on a data-rich task before being fine-tuned on a downstream task, has emerged as a powerful technique in NLP. In this paper, we explore the landscape of transfer learning techniques for NLP by introducing a unified framework that converts all text-based language problems into a text-to-text format.

\vspace{18pt}

% ===================================================================
% CITATION KEY: [vaswani2017attention]
% Length: 1920 characters
% ===================================================================

\subsection*{Paper Reference: vaswani2017attention}

1 Introduction Recurrent neural networks, long short-term memory and gated recurrent neural networks in particular, have been firmly established as state of the art approaches in sequence modeling and transduction problems such as language modeling and machine translation. Numerous efforts have since continued to push the boundaries of recurrent language models and encoder-decoder architectures. In this work we propose the Transformer, a model architecture eschewing recurrence and instead relying entirely on an attention mechanism to draw global dependencies between input and output.

\vspace{18pt}

% ===================================================================
% CITATION KEY: [xue2021mt5]
% Length: 3228 characters
% ===================================================================

\subsection*{Paper Reference: xue2021mt5}

1 Introduction Current natural language processing (NLP) pipelines often make use of transfer learning, where a model is pre-trained on a data-rich task before being fine-tuned on a downstream task of interest. The recent Text-to-Text Transfer Transformer (T5) leveraged a unified text-to-text format and scale to attain state-of-the-art results on a wide variety of English-language NLP tasks. In this paper, we introduce mT5, a multilingual variant of T5 that was pre-trained on a new Common Crawl-based dataset covering 101 languages.

\end{document}
